% preamble.tex — Shared LaTeX preamble for textbook conversion
% Page geometry configured to match the source textbook

\documentclass[10pt,openright]{book}

%% ── Page Geometry ──────────────────────────────────────────────
% Original is approximately 6.5" x 9.5" Prentice Hall format
\usepackage[
  paperwidth=6.5in,
  paperheight=9.5in,
  top=1in,
  bottom=1in,
  left=0.85in,
  right=0.85in,
  headheight=14pt,
  headsep=0.3in,
  footskip=0.4in
]{geometry}

%% ── Encoding & Fonts ──────────────────────────────────────────
\usepackage[T1]{fontenc}
\usepackage[utf8]{inputenc}
\usepackage{lmodern}          % Latin Modern (close to Computer Modern)
\usepackage{microtype}        % Improved typography

%% ── Math ───────────────────────────────────────────────────────
\usepackage{amsmath}
\usepackage{amssymb}
\usepackage{amsthm}
\usepackage{mathtools}        % Extensions to amsmath
\usepackage{bm}               % Bold math symbols

% Equation numbering: (chapter.section.number)
\numberwithin{equation}{section}

% Figure/table numbering: chapter.section.number
\usepackage{chngcntr}
\counterwithin{figure}{section}
\counterwithin{table}{section}

%% ── Boxed Equations/Definitions ────────────────────────────────
\usepackage[most]{tcolorbox}

% Simple box style matching the book's thin-border boxes
\newtcolorbox{defbox}{
  colback=white,
  colframe=black,
  boxrule=0.5pt,
  arc=0pt,
  left=6pt, right=6pt, top=4pt, bottom=4pt,
  boxsep=2pt
}

%% ── Theorem Environments ───────────────────────────────────────
\theoremstyle{plain}
\newtheorem{theorem}{Theorem}[section]
\newtheorem{lemma}[theorem]{Lemma}
\newtheorem{corollary}[theorem]{Corollary}
\newtheorem{proposition}[theorem]{Proposition}

\theoremstyle{definition}
\newtheorem{definition}{Definition}[section]
\newtheorem{example}{Example}[section]

\theoremstyle{remark}
\newtheorem{remark}{Remark}[section]

%% ── Exercises ──────────────────────────────────────────────────
\usepackage{enumitem}

% Exercise label style — configure per book, then forget about it.
%   \exercisestyle{prefix}{counter} is called for each item label.
%   Default: "1.2.1." "1.2.2." (section-prefixed)
%   Override examples:
%     Plain numbers:       \renewcommand{\exerciselabel}[1]{\textbf{\arabic*.}}
%     Parenthesized:       \renewcommand{\exerciselabel}[1]{\textbf{(\arabic*)}}
%     Chapter-only prefix: \renewcommand{\exerciselabel}[1]{\textbf{\thechapter.\arabic*.}}
\newcommand{\exerciselabel}[1]{\textbf{#1.\arabic*.}}
\newcommand{\exerciseref}[1]{#1.\arabic*}

\newenvironment{exercises}[1]{%
  \subsection*{Exercises #1}%
  \begin{enumerate}[label={\exerciselabel{#1}},ref={\exerciseref{#1}},leftmargin=*]%
}{%
  \end{enumerate}%
}

% Starred exercise marker
\newcommand{\starred}{\textsuperscript{*}}

%% ── Figures ────────────────────────────────────────────────────
\usepackage{graphicx}
\usepackage{float}
\usepackage{tikz}
\usepackage{pgfplots}
\pgfplotsset{compat=1.18}
\usetikzlibrary{arrows.meta, patterns, decorations.markings, decorations.pathmorphing, calc}

%% ── Tables ─────────────────────────────────────────────────────
\usepackage{booktabs}
\usepackage{array}
\usepackage{multirow}

%% ── Headers & Footers ──────────────────────────────────────────
\usepackage{fancyhdr}
\pagestyle{fancy}
\fancyhf{}
\fancyhead[LE]{\thepage}
\fancyhead[RE]{\leftmark}
\fancyhead[RO]{\thepage}
\fancyhead[LO]{\rightmark}
\renewcommand{\headrulewidth}{0pt}

% Chapter/section mark formatting
\renewcommand{\chaptermark}[1]{\markboth{Chapter \thechapter.\ #1}{}}
\renewcommand{\sectionmark}[1]{\markright{\thesection\ #1}}

% Plain style for chapter opening pages
\fancypagestyle{plain}{%
  \fancyhf{}%
  \fancyfoot[C]{\thepage}%
  \renewcommand{\headrulewidth}{0pt}%
}

%% ── Cross-references ──────────────────────────────────────────
\usepackage{hyperref}
\hypersetup{
  colorlinks=false,
  pdfborder={0 0 0},
  % PDF metadata — set per book
  pdftitle={},
  pdfauthor={}
}
\usepackage{cleveref}

%% ── Index ──────────────────────────────────────────────────────
\usepackage{makeidx}
\makeindex

%% ── Custom Commands ────────────────────────────────────────────

% Partial derivatives
\newcommand{\pd}[2]{\frac{\partial #1}{\partial #2}}
\newcommand{\pdd}[2]{\frac{\partial^2 #1}{\partial #2^2}}
\newcommand{\pddm}[3]{\frac{\partial^2 #1}{\partial #2 \partial #3}}

% Ordinary derivatives
\newcommand{\od}[2]{\frac{d #1}{d #2}}
\newcommand{\odd}[2]{\frac{d^2 #1}{d #2^2}}

% Evaluation bar
\newcommand{\evalat}[2]{\left. #1 \right|_{#2}}

% Common operators
\DeclareMathOperator{\erfc}{erfc}
\DeclareMathOperator{\sinc}{sinc}
\DeclareMathOperator{\sgn}{sgn}
\DeclareMathOperator{\Real}{Re}
\DeclareMathOperator{\Imag}{Im}

% Bold vectors
\renewcommand{\vec}[1]{\mathbf{#1}}

% Laplacian
\newcommand{\laplacian}{\nabla^2}

% Double surface integral (not in all packages)
\providecommand{\oiint}{\mathop{\oint\!\!\!\oint}\nolimits}

% Common notations used in the book
\newcommand{\dx}{\,dx}
\newcommand{\dy}{\,dy}
\newcommand{\dt}{\,dt}
\newcommand{\ds}{\,ds}
\newcommand{\dA}{\,dA}
\newcommand{\dV}{\,dV}
\newcommand{\dS}{\,dS}
